% Nested rectangles distinguish containment from the size of infinite sets.
\begin{tikzpicture}[analysis diagram]
  \draw (0,0) rectangle (10,7.2);
  \node[analysis label,anchor=north west] at (0.3,6.95)
    {$\C$: complex numbers};
  \node[analysis label,anchor=north east] at (9.65,6.95) {$i,\ 1+i$};

  \draw (0.35,0.35) rectangle (9.65,6.25);
  \node[analysis label,anchor=north west] at (0.65,6)
    {$\R$: real numbers};

  \draw (0.7,0.7) rectangle (6.3,5.3);
  \node[analysis label,anchor=north west] at (1,5)
    {$\mathcal A$: real algebraic numbers};
  \node[analysis label,anchor=north west] at (1,4.35) {$\sqrt2,\ \sqrt[3]{5}$};

  \draw (1.05,1.05) rectangle (5.95,3.7);
  \node[analysis label,anchor=north west] at (1.35,3.45)
    {$\Q$: rational numbers};
  \node[analysis label,anchor=east] at (5.6,2.5) {$\frac12,\ -\frac34$};

  \draw (1.4,1.4) rectangle (4.25,2.8);
  \node[analysis label,anchor=north west] at (1.7,2.55)
    {$\N$: naturals};
  \node[analysis label,anchor=north west] at (1.7,1.95)
    {$1,2,3,\ldots$};

  \node[analysis label,align=center] at (7.95,4.65)
    {Transcendental reals\\$\mathcal T=\R\setminus\mathcal A$};
  \node[analysis label] at (7.95,3.6) {$\pi,\ e$};
\end{tikzpicture}
